\section{Correspondence with VED}
\label{sec:correspondence}

The division of roles described in the previous section is now organized in terms of variable correspondence and dynamical correspondence.

\subsection{Variable correspondence table}
\label{subsec:variable_table}

\begin{table}[h]
  \centering
  \begin{tabular}{l l l}
    \hline
    \textbf{VED} & \textbf{IFGT} & \textbf{Relation} \\
    \hline
    closure degree $C \in [0,1]$ & information density $I$ & $I = f(1 - C)$ \\
    difference $\Delta$ & information difference $\delta I$ & difference residue as trace \\
    gradient $\nabla C$ & information gradient $\nabla I$ & opposite direction via $I = f(1-C)$ \\
    closure flow $J_C$ & information flow $J$ & different aspects of the same flow \\
    closure potential & information potential $\Phi$ & $\Phi = -\alpha C + \Phi_0$ \\
    generative time $\tau_{\mathrm{VED}}$ & informational time $\tau$ & different origin, dynamically isomorphic \\
    --- & generation term $\sigma$ & corresponds to VED's difference generation \\
    \hline
  \end{tabular}
  \caption{Variable correspondence between VED and IFGT.}
  \label{tab:variable_correspondence}
\end{table}

\subsection{Difference \texorpdfstring{$\Delta$}{Delta} and information difference \texorpdfstring{$\delta I$}{delta I}}
\label{subsec:difference}

In VED, the difference $\Delta$ is a primitive quantity derived directly from the foundational axiom ``\emph{there is difference}'' and serves as the driver of closure formation. In IFGT, the information difference $\delta I$ is the residual difference that remains when difference does not fully close.

Whereas $\Delta$ is unprocessed difference itself, $\delta I$ is interpreted as ``the part that did not contribute to closure.'' The following relation therefore holds:
\begin{equation}
\Delta = \Delta_{\to C} + \delta I,
\label{eq:delta_decomposition}
\end{equation}
where $\Delta_{\to C}$ is the component of difference that contributes to closure and $\delta I$ is the component that persists as a quasi-closure.

\subsection{Structural correspondence}
\label{subsec:structural_correspondence}

In VED, difference produces gradient, gradient produces flow, and flow forms closure, giving rise to structure as physical persistence.

In IFGT, when difference does not fully close, it persists as a quasi-closure. This quasi-closure accumulates as a trace, constituting the information density $I$. The non-uniformity of traces produces the information gradient $\nabla I$, which drives the information flow $J$. Furthermore, the accumulation of quasi-closures constrains the flow, which is described as the information potential $\Phi$.

\subsection{Dynamical correspondence}
\label{subsec:dynamical_correspondence}

\begin{table}[h]
  \centering
  \begin{tabular}{l l l}
    \hline
    \textbf{Aspect} & \textbf{VED} & \textbf{IFGT} \\
    \hline
    structure & closure (stable, tending to rest) & quasi-closure (circulatory, persisting motion) \\
    terminal state & complete closure $C \to 1$ & metastable circulation \\
    time evolution & closure formation & trace update \\
    \hline
  \end{tabular}
  \caption{Dynamical correspondence between VED and IFGT.}
  \label{tab:dynamical_correspondence}
\end{table}

Closure is a stable structure tending toward rest, whereas a quasi-closure forms a circulatory structure while maintaining flow. This circulation strengthens local constraints and draws in surrounding flow.

\subsection{Generation and mismatch}
\label{subsec:generation_mismatch}

The generation term $\sigma$ in IFGT represents the mismatch between the current difference chain and the existing quasi-closure structure. This mismatch corresponds to the generation of new differences in VED and induces the update of structure.

Thus the generation of difference in VED and the generation through $\sigma$ in IFGT are different descriptions of the same process.
